\subsection{Theoretical intuition for mode collapse}
\label{sec:collapse-mechanism}

Binary rewards treat all correct modes alike, while fresh samples follow the policy's current preferences. A mode with an early probability advantage is therefore sampled and reinforced more often (Fig.~\ref{fig:replay-bank-balance}). Under independent logit updates, this feedback can produce winner-take-all concentration. Changing the binary-reward estimator alone does not change the trajectory of this concentration, as the following result makes precise.

For one isolated prompt, let a finite categorical policy have $m\ge2$ correct modes and $n\ge1$ incorrect categories, with $p=\operatorname{softmax}(z)$, correctness $P=\sum_c p_c$, and conditional mode probabilities $q_c=p_c/P$. Each category has an independent logit. Draw responses $Z_1,\ldots,Z_G$ independently with $G\ge2$, with binary rewards $R_i$ and $R=\sum_iR_i$. Assume finite initial logits, common response-length normalization, and on-policy Euclidean expected-gradient flow without KL, entropy, momentum, or preconditioning (Assumption~\ref{ass:categorical-mean-flow}).

\begin{theorem}[Conditional inertness of outcome-binary objectives]
\label{thm:objective-inertness}
Under Assumption~\ref{ass:categorical-mean-flow}, let $a(\cdot,\cdot)$ be any
detached advantage measurable with respect to the group's binary reward
vector, applied as
$\widehat g_a=G^{-1}\sum_{i}a(R_i,R)\nabla_z\log p_{Z_i}$ at the on-policy
point. Then
\begin{equation}
  \E[\widehat g_a]=c_a(P)\,\nabla_zP,
  \qquad
  c_a(P)=\frac1G\left(
  \frac{\E[R\,a(1,R)]}{P}-\frac{\E[(G-R)\,a(0,R)]}{1-P}\right).
  \label{eq:objective-inertness}
\end{equation}
If $c_a>0$ on $(0,1)$, then under the time change
$d\tau=c_a(P)P(1-P)\,dt$ the conditional distribution $q$ obeys
Equation~\eqref{eq:grpo-replicator} exactly, for every such $a$. Dr.GRPO,
reward-standardized GRPO, a leave-one-out baseline, the centered binary MaxRL
estimator, and any reweighting shaped by a pass@$k$ target are points in this
class. The choice among them rescales optimization time along the correctness
axis and leaves both the trajectory in $q$ and the identity of the surviving
mode exactly where they were.
\end{theorem}

The conditional dynamics are $dq_c/d\tau=q_c(q_c-\sum_dq_d^2)$: above-average modes grow at the expense of others. With a unique initial maximum and unbounded $\tau$, one mode eventually dominates (Theorem~\ref{thm:grpo-collapse}). The proof of Theorem~\ref{thm:objective-inertness} is in App.~\ref{app:theory-objective-inertness}. These are statements about the specified logit geometry, not guarantees for stochastic neural training. Replay changes the available targets by retaining successes that fresh sampling may stop producing.

